% GENERATED FILE. Edit canonical lessons, then run npm run generate:books.
% canonical-source-hash: fnv1a64:8c542451978748ec
% canonical-lessons: LA-C393-statuo, LA-C393-destino, LA-C393-propositum, LA-C393-inceptum, LA-C393-conor

\chapter{Making Up My Mind}
\label{ch:la-making-up-my-mind}
\hlchaptermodality{393}

\begin{chapteropening}
\textbf{By the end of this chapter:} \emph{I can say that I have decided and fixed a day, talk about an aim and a venture, and say that I am trying to do something.}

Complete the last lesson of chapter 393: I can say that I have decided and fixed a day, talk about an aim and a venture, and say that I am trying to do something.
\end{chapteropening}

\section[statuō, statuere]{statuō, statuere — to set up; to decide}

\label{lesson:LA-C393-statuo}



\begin{quote}
Before the new one: say the Latin for convenient, then the Latin for to take precautions.
\end{quote}

\subsection*{You'll want to know: statuō, statuere}
\textbf{statuō, statuere} — \textquotedblleft{}to set up; to decide\textquotedblright{}.

\begin{quote}
Statuī domī manēre.
\end{quote}

\begin{quote}
\emph{I decided to stay at home.}
\end{quote}

\emph{statuō, statuere, statuī, statūtum}. With an infinitive, to decide to; with the accusative and infinitive, to judge that. English \emph{statute}.

\begin{cousinweb}
\emph{statuō} is built on \emph{status} \textquotedblleft{}a standing, a position\textquotedblright{} (Lewis \& Short), from the root of \emph{stō} \textquotedblleft{}to stand\textquotedblright{}: to make a thing stand. English \emph{statute} comes from it.
\end{cousinweb}

\begin{grammarlens}[title={What you've built}]
One more word for telling, explaining and describing at B1.
\end{grammarlens}

\subsection*{Your turn}
\begin{itemize}
\raggedright
  \item \emph{Say it:} \emph{statuō, statuere}
  \item \emph{Say it:} \emph{statuō, statuere}, once more
  \item \emph{Say it:} a sentence of your own with \emph{statuō, statuere}
  \item \emph{Recall:} say the Latin for convenient, then the Latin for to take precautions, then say \emph{statuō, statuere} again
\end{itemize}

\begin{tcolorbox}[breakable,colback=teal!4,colframe=teal!35!black,title={Before you move on}]
What does statuō, statuere mean? (\textquotedblleft{}To set up; to decide\textquotedblright{}.) Say it once more.
\end{tcolorbox}
\section[dēstinō, dēstināre]{dēstinō, dēstināre — to fix, to intend; to appoint}

\label{lesson:LA-C393-destino}



\begin{quote}
Before the new one: say the Latin for to take precautions, then the Latin for to set up.
\end{quote}

\subsection*{You'll want to know: dēstinō, dēstināre}
\textbf{dēstinō, dēstināre} — \textquotedblleft{}to fix, to intend; to appoint\textquotedblright{}.

\begin{quote}
Hunc diem nūptiīs dēstināvērunt.
\end{quote}

\begin{quote}
\emph{They fixed this day for the wedding.}
\end{quote}

\emph{dēstinō, dēstināre, dēstināvī, dēstinātum}: to make fast, then to set aside for a purpose. English \emph{destine} and \emph{destination}.

\begin{cousinweb}
\emph{dēstinō} is built from \emph{dē} and an old verb meaning \textquotedblleft{}to set, to make stand\textquotedblright{} (EtymDB), from the root of \emph{stō} \textquotedblleft{}to stand\textquotedblright{}. English \emph{destine} and \emph{destination} come from it.
\end{cousinweb}

\begin{grammarlens}[title={What you've built}]
One more word for telling, explaining and describing at B1.
\end{grammarlens}

\subsection*{Your turn}
\begin{itemize}
\raggedright
  \item \emph{Say it:} \emph{dēstinō, dēstināre}
  \item \emph{Say it:} \emph{dēstinō, dēstināre}, once more
  \item \emph{Say it:} a sentence of your own with \emph{dēstinō, dēstināre}
  \item \emph{Recall:} say the Latin for to take precautions, then the Latin for to set up, then say \emph{dēstinō, dēstināre} again
\end{itemize}

\begin{tcolorbox}[breakable,colback=teal!4,colframe=teal!35!black,title={Before you move on}]
What does dēstinō, dēstināre mean? (\textquotedblleft{}To fix, to intend; to appoint\textquotedblright{}.) Say it once more.
\end{tcolorbox}
\section[prōpositum, prōpositī]{prōpositum, prōpositī — a plan, a purpose, an aim}

\label{lesson:LA-C393-propositum}



\begin{quote}
Before the new one: say the Latin for to set up, then the Latin for to fix.
\end{quote}

\subsection*{You'll want to know: prōpositum, prōpositī}
\textbf{prōpositum, prōpositī} — \textquotedblleft{}a plan, a purpose, an aim\textquotedblright{}.

\begin{quote}
Prōpositum meum est librum scrībere.
\end{quote}

\begin{quote}
\emph{My aim is to write a book.}
\end{quote}

\emph{prōpositum, prōpositī} n., the participle of \emph{prōpōnō}, to set before: what you have set before yourself. English \emph{purpose} goes back to it through French.

\begin{cousinweb}
\emph{prōpositum} is the participle of \emph{prōpōnō} \textquotedblleft{}to set before\textquotedblright{} (Lewis \& Short), from \emph{prō} \textquotedblleft{}before\textquotedblright{} and \emph{pōnō} \textquotedblleft{}to put\textquotedblright{}. English \emph{purpose} goes back to the same verb through French.
\end{cousinweb}

\begin{grammarlens}[title={What you've built}]
One more word for telling, explaining and describing at B1.
\end{grammarlens}

\subsection*{Your turn}
\begin{itemize}
\raggedright
  \item \emph{Say it:} \emph{prōpositum, prōpositī}
  \item \emph{Say it:} \emph{prōpositum, prōpositī}, once more
  \item \emph{Say it:} a sentence of your own with \emph{prōpositum, prōpositī}
  \item \emph{Recall:} say the Latin for to set up, then the Latin for to fix, then say \emph{prōpositum, prōpositī} again
\end{itemize}

\begin{tcolorbox}[breakable,colback=teal!4,colframe=teal!35!black,title={Before you move on}]
What does prōpositum, prōpositī mean? (\textquotedblleft{}A plan, a purpose, an aim\textquotedblright{}.) Say it once more.
\end{tcolorbox}
\section[inceptum, inceptī]{inceptum, inceptī — an undertaking, a venture}

\label{lesson:LA-C393-inceptum}



\begin{quote}
Before the new one: say the Latin for to fix, then the Latin for a plan.
\end{quote}

\subsection*{You'll want to know: inceptum, inceptī}
\textbf{inceptum, inceptī} — \textquotedblleft{}an undertaking, a venture\textquotedblright{}.

\begin{quote}
Inceptum tuum laudō.
\end{quote}

\begin{quote}
\emph{I praise your undertaking.}
\end{quote}

\emph{inceptum, inceptī} n., the participle of \emph{incipiō}, to begin: a thing begun.

\begin{cousinweb}
\emph{inceptum} is the participle of \emph{incipiō} \textquotedblleft{}to begin\textquotedblright{} (Lewis \& Short), from \emph{in} and \emph{capiō} \textquotedblleft{}to take\textquotedblright{}: a thing taken in hand. English \emph{inception} comes from its family.
\end{cousinweb}

\begin{grammarlens}[title={What you've built}]
One more word for telling, explaining and describing at B1.
\end{grammarlens}

\subsection*{Your turn}
\begin{itemize}
\raggedright
  \item \emph{Say it:} \emph{inceptum, inceptī}
  \item \emph{Say it:} \emph{inceptum, inceptī}, once more
  \item \emph{Say it:} a sentence of your own with \emph{inceptum, inceptī}
  \item \emph{Recall:} say the Latin for to fix, then the Latin for a plan, then say \emph{inceptum, inceptī} again
\end{itemize}

\begin{tcolorbox}[breakable,colback=teal!4,colframe=teal!35!black,title={Before you move on}]
What does inceptum, inceptī mean? (\textquotedblleft{}An undertaking, a venture\textquotedblright{}.) Say it once more.
\end{tcolorbox}
\section[cōnor, cōnārī]{cōnor, cōnārī — to try, to attempt}

\label{lesson:LA-C393-conor}



\begin{quote}
Before the new one: say the Latin for a plan, then the Latin for an undertaking.
\end{quote}

\subsection*{You'll want to know: cōnor, cōnārī}
\textbf{cōnor, cōnārī} — \textquotedblleft{}to try, to attempt\textquotedblright{}.

\begin{quote}
Cōnābar dormīre, sed nōn poteram.
\end{quote}

\begin{quote}
\emph{I was trying to sleep, but I could not.}
\end{quote}

Deponent: \emph{cōnor, cōnārī, cōnātus sum} + infinitive. \emph{Cōnātus}, an attempt.

\begin{cousinweb}
\emph{cōnor} is inherited from Proto-Italic \emph{*kōnāō} \textquotedblleft{}to try, to exert oneself\textquotedblright{} (EtymDB); its deeper origin is not settled.
\end{cousinweb}

\begin{grammarlens}[title={What you've built}]
One more word for telling, explaining and describing at B1.
\end{grammarlens}

\subsection*{Your turn}
\begin{itemize}
\raggedright
  \item \emph{Say it:} \emph{cōnor, cōnārī}
  \item \emph{Say it:} \emph{cōnor, cōnārī}, once more
  \item \emph{Say it:} a sentence of your own with \emph{cōnor, cōnārī}
  \item \emph{Recall:} say the Latin for a plan, then the Latin for an undertaking, then say \emph{cōnor, cōnārī} again
\end{itemize}

\begin{tcolorbox}[breakable,colback=teal!4,colframe=teal!35!black,title={Before you move on}]
What does cōnor, cōnārī mean? (\textquotedblleft{}To try, to attempt\textquotedblright{}.) Say it once more.
\end{tcolorbox}
